\subsection{Elektronen im elektromagnetischen Feld: Die Lorentzkraft}
\label{sec:elektronen-lorentzkraft}

Elektronen sind geladene Teilchen. Deshalb reagieren sie auf elektrische
und magnetische Felder. In der klassischen Physik wird ihre Bewegung durch
eine Kraftgleichung beschrieben, die Mechanik und Elektrodynamik miteinander
verbindet: die Lorentzkraft\index{Lorentzkraft}.

Ein elektrisches Feld\index{Elektrisches Feld} übt auf eine Ladung \(q\)
die Kraft

\[
\vec{F}_E = q \vec{E}
\]

aus. Für ein Elektron gilt \(q=-e\). Die Kraft zeigt deshalb entgegen der
Richtung des elektrischen Feldes:

\[
\vec{F}_E = -e \vec{E}.
\]

Das ist ein wichtiger Punkt: Die Feldlinien eines elektrischen Feldes geben
die Richtung der Kraft auf eine positive Probeladung an. Ein Elektron wird
wegen seiner negativen Ladung in die entgegengesetzte Richtung beschleunigt.

Befindet sich das Elektron dagegen in einem magnetischen Feld\index{Magnetisches Feld},
so hängt die Kraft nicht nur von der Ladung und vom Magnetfeld ab, sondern
auch von der Geschwindigkeit des Elektrons. Für eine bewegte Ladung gilt:

\[
\vec{F}_B = q \, \vec{v} \times \vec{B}.
\]

Für das Elektron wird daraus

\[
\vec{F}_B = -e \, \vec{v} \times \vec{B}.
\]
Eine ausführliche Herleitung dieser Beziehung findet sich in
\hyperref[app:herleitung-magnetische-lorentzkraft]{Anhang~\ref*{app:herleitung-magnetische-lorentzkraft}}.
\FloatBarrier
\begin{figure}[H]
	\centering
	\includegraphics[width=\textwidth]{bilder/lorentzkraft.png}
	\caption{Richtung der magnetischen Lorentzkraft auf ein Elektron. Das Elektron bewegt sich mit der Geschwindigkeit \(\vec{v}\) nach rechts, während das Magnetfeld \(\vec{B}\) in die Bildebene hinein zeigt. Für eine positive Ladung würde die Kraft in Richtung \(\vec{v} \times \vec{B}\) zeigen. Da das Elektron negativ geladen ist, wirkt die magnetische Kraft in die entgegengesetzte Richtung.}
	\label{fig:lorentzkraft-richtung-elektron}
\end{figure}
Dieses Bild ist besonders hilfreich, weil es den entscheidenden Punkt
sichtbar macht: Das Magnetfeld zieht das Elektron nicht zu einem Pol, sondern
lenkt seine Bewegung seitlich ab. Die Kraft steht stets senkrecht zur
momentanen Geschwindigkeit.
Diese Kraft steht senkrecht zur Bewegungsrichtung \(\vec{v}\) und senkrecht
zum Magnetfeld \(\vec{B}\). Sie verändert daher nicht direkt den Betrag der
Geschwindigkeit, sondern vor allem die Richtung der Bewegung. Genau deshalb
können Elektronen in Magnetfeldern auf Kreisbahnen oder Schraubenbahnen
gezwungen werden.

Elektrisches und magnetisches Feld zusammen ergeben die Lorentzkraft:

\[
\vec{F}
=
q\left(\vec{E} + \vec{v} \times \vec{B}\right).
\]

Für das Elektron lautet sie:

\[
\vec{F}
=
-e\left(\vec{E} + \vec{v} \times \vec{B}\right).
\]

Setzt man diese Kraft in das zweite Newtonsche Gesetz ein, erhält man die
Bewegungsgleichung des Elektrons im elektromagnetischen Feld:

\[
m_e \vec{a}
=
-e\left(\vec{E} + \vec{v} \times \vec{B}\right).
\]

Diese Gleichung ist einer der wichtigsten Ausgangspunkte der klassischen
Elektronenphysik. Sie beschreibt die Ablenkung von Elektronenstrahlen in
Kathodenstrahlröhren, die Bewegung von Elektronen in Magnetfeldern, die
Funktionsweise vieler Messgeräte und die Grundlage zahlreicher technischer
Anwendungen.

\begin{PhysicsBox}[Die Lorentzkraft]
	\label{box:lorentzkraft}
	Die Lorentzkraft beschreibt die Kraft auf eine bewegte elektrische Ladung
	in elektrischen und magnetischen Feldern:
	
	\[
	\vec{F}
	=
	q\left(\vec{E} + \vec{v} \times \vec{B}\right).
	\]
	
	Für ein Elektron mit \(q=-e\) gilt:
	
	\[
	\vec{F}
	=
	-e\left(\vec{E} + \vec{v} \times \vec{B}\right).
	\]
	
	Der elektrische Anteil kann das Elektron beschleunigen oder abbremsen.
	Der magnetische Anteil lenkt die Bewegung seitlich ab und steht senkrecht
	zur Geschwindigkeit.
\end{PhysicsBox}

Anschaulich kann man sich die Wirkung so merken:
Ein elektrisches Feld zieht oder schiebt das Elektron entlang einer Linie.
Ein magnetisches Feld dagegen wirkt nur auf bewegte Elektronen und lenkt
ihre Bahn seitlich ab. Das elektrische Feld ändert also typischerweise die
Geschwindigkeit, das magnetische Feld ändert vor allem die Richtung.

Diese Unterscheidung ist entscheidend. Sie erklärt, warum Elektronen in
elektrischen Feldern beschleunigt werden können, während Magnetfelder sie
auf gekrümmte Bahnen zwingen. Im nächsten Abschnitt wird daraus unmittelbar
die Kreisbewegung eines Elektrons im Magnetfeld folgen.